% Source: ACM_Templates.pdf, PDF pp. 93-95; generalized from task-specific counters.
\subsection{树上启发式合并：DSU on tree}
\label{tree:dsu-on-tree}
\AlgorithmInfo
    {离线求每个子树的统计信息，示例为子树内不同颜色数。}
    {每次加点、删点为 $O(1)$ 时，总时间 $O(n\log n)$，空间 $O(n+m)$。}
    {先复制 HLD；颜色为 $1\ldots m$ 的整数，数组从 1 开始，原树连通。}
\subsubsection*{原理}
按子树大小选重儿子。递归处理所有轻儿子并清空它们的贡献，接着处理重儿子并保留其贡献；再逐个加入轻子树和当前点，此时容器恰好包含当前整棵子树，可以记录答案。若父层不希望保留当前子树，则在返回前统一删除它。

每个点只有经过轻边时才被重复加入或删除，轻边经过次数为 $O(\log n)$，因此在单点增删常数时间的前提下总复杂度为 $O(n\log n)$。它与普通并查集的合并不是同一种接口，也不能把这里按子树大小选的重儿子改成按高度选择的长儿子。

模板借用 HLD 的父亲、重儿子和子树区间，递归控制保留顺序；按 DFS 序扫描子树来增删颜色计数。\InlineCode{eval} 维护频次数组与当前不同颜色数，只在完整子树合并后记录 \InlineCode{ans[u]}。
\subsubsection*{使用教程}
\InlineCode{init(n,m)} 指定顶点数和颜色上界，再用 \InlineCode{add(u,v)} 加无向边，每边一次。\InlineCode{run(color,root)} 接收第 0 项占位的颜色数组，返回从 1 开始的子树不同颜色数；省略根取 1，每次运行重建剖分并清空统计。颜色须在 $1\ldots m$ 内，范围大时先离散化。改换统计时，同步调整 \InlineCode{eval}、清理和记录答案的位置，保持增删互逆。
\begin{lstlisting}[style=mtycpp,numbers=none]
DSUOnTree g;
g.init(5,3);
g.add(1,2),g.add(1,3),g.add(2,4),g.add(2,5);
vector<int> color={0,1,2,1,3,2};
auto ans=g.run(color,1);
int distinct=ans[2];
\end{lstlisting}
\needspace{7\baselineskip}
\subsubsection*{模板代码}
\CodeFile{code/trees/dsu_on_tree.cpp}
